% figures/ch05-hecke-q-expansion.tex
% Hecke 算子 T_p 在 q-展开上的作用示意
\begin{figure}[ht]
\centering
\begin{tikzpicture}[>=Stealth, font=\small]

% 标题
\node[font=\normalsize\bfseries] at (0, 3.5) {$T_p$ 在 $q$-展开系数上的作用};

% f 的系数序列
\node[align=left] at (-4.5, 2.7) {$f = \sum a_n q^n$:};
\foreach \n/\x in {0/-3, 1/-1.5, 2/0, 3/1.5, 4/3, 5/4.5}{
  \draw[fill=blue!15, draw=blue!50!black, rounded corners=2pt]
    (\x-0.55, 2.3) rectangle (\x+0.55, 3.0);
  \node[font=\footnotesize] at (\x, 2.65) {$a_{\n}$};
  \node[font=\footnotesize, below] at (\x, 2.3) {$q^{\n}$};
}
\node at (5.5, 2.65) {$\cdots$};

% 箭头
\draw[->, very thick, purple!70!black] (0, 2.0) -- (0, 1.3)
  node[right, midway] {$T_p$};

% T_p f 的系数
\node[align=left] at (-4.5, 0.9) {$T_p f = \sum b_n q^n$:};

\foreach \n/\x/\formula in {
  0/-3/{$a_{p}$},
  1/-1.5/{$a_{p{+}1}$},
  2/0/{$a_{2p}$},
  3/1.5/{$a_{3p}$},
  4/3/{$a_{4p}$},
  5/4.5/{$a_{5p}$}
}{
  \draw[fill=orange!20, draw=orange!60!black, rounded corners=2pt]
    (\x-0.65, 0.4) rectangle (\x+0.65, 1.15);
  \node[font=\footnotesize] at (\x, 0.8) {\formula};
  \node[font=\footnotesize, below] at (\x, 0.4) {$q^{\n}$};
}
\node at (5.5, 0.8) {$\cdots$};

% 完整公式
\node[draw=gray!50, rounded corners=4pt, fill=gray!5, font=\footnotesize, align=center]
  at (0, -0.6) {
  $b_n = a_{pn} + p^{k-1} a_{n/p}$
  \quad（其中 $a_{n/p} = 0$ 若 $p \nmid n$）
};

% 特殊情形提示
\node[font=\footnotesize, text=green!50!black] at (0, -1.3)
  {特别：$b_1 = a_p$（$T_p$ 的特征值就是 $a_p$，若 $f$ 是特征形式）};

\end{tikzpicture}
\caption{Hecke 算子 $T_p$ 在权 $k$ 模形式 $f = \sum_{n=0}^\infty a_n q^n$ 的
$q$-展开系数上的作用公式：$(T_p f)$ 的第 $n$ 个系数为 $a_{pn} + p^{k-1} a_{n/p}$。
注意 $T_p f$ 的第 $1$ 个系数是 $a_p$——若 $f$ 是特征形式，
则 $T_p f = a_p f$，所以 $a_p$ 就是 $T_p$ 的特征值。}
\label{fig:hecke-q-expansion}
\end{figure}
